\documentclass[conference]{IEEEtran}
\usepackage[utf8]{inputenc}
\usepackage[T1]{fontenc}
\usepackage[brazil]{babel}
\usepackage{booktabs}
\usepackage{graphicx}
\usepackage{url}
\usepackage{listings}

% ===================================================================
% INSTRUÇÕES: os trechos marcados com [PREENCHER] dependem dos seus
% resultados experimentais. Substitua-os após executar o pipeline.
% ===================================================================

\title{Análise Quantitativa do Trade-off entre Especialização e
Generalização em LLMs via Fine-Tuning}

\author{\IEEEauthorblockN{[PREENCHER: Seu Nome]}
\IEEEauthorblockA{Instituto de Computação -- Universidade Federal do Amazonas\\
[PREENCHER: seu e-mail]}}

\begin{document}
\maketitle

\begin{abstract}
Este trabalho avalia empiricamente o trade-off entre especialização e
generalização no fine-tuning de Modelos de Linguagem de Grande Porte.
Um modelo Qwen2.5-3B-Instruct foi especializado para a tarefa de
Text-to-SQL via QLoRA sobre o training split do benchmark Spider, em
duas configurações de hiperparâmetros, e avaliado com uma métrica
customizada de Execution Accuracy implementada em DeepEval. A regressão
de capacidade geral foi medida em uma suíte de 150 questões do MMLU
(5-shot) cobrindo STEM, Humanidades e Ciências Sociais. Os resultados
mostram um ganho de [PREENCHER] pontos percentuais na tarefa-alvo,
acompanhado de uma variação de [PREENCHER]\% na acurácia agregada do
MMLU, com efeitos heterogêneos entre categorias.
\end{abstract}

\section{Introdução}
A especialização de LLMs via fine-tuning é amplamente empregada para
otimizar o desempenho em domínios específicos, mas pode comprometer a
robustez do modelo fora do domínio de treinamento --- fenômeno conhecido
como esquecimento catastrófico ou regressão de capacidade. Este trabalho
projeta e executa um pipeline experimental que mede com precisão ambas
as facetas do fenômeno: o ganho de especialização em Text-to-SQL e a
perda de generalização em conhecimento geral.

\section{Metodologia}

\subsection{Modelo e infraestrutura}
Utilizamos o modelo \texttt{Qwen/Qwen2.5-3B-Instruct} como base. Todos
os experimentos foram executados em [PREENCHER: GPU utilizada, e.g.,
NVIDIA T4 com 16 GB de VRAM, Google Colab tier gratuito].

\subsection{Pipeline de dados}
O training split do Spider (7.000 exemplos) foi convertido para o
formato de chat com três turnos: uma mensagem de sistema fixa instruindo
a geração de SQL em bloco markdown, uma mensagem de usuário contendo a
serialização textual do esquema do banco-alvo seguida da pergunta em
linguagem natural, e a resposta de referência como turno do assistente.
A serialização do esquema lista, para cada tabela, suas colunas com
tipos, seguidas das chaves primárias e estrangeiras, suprindo a
informação necessária ao \emph{schema linking}. O dev split (1.034
exemplos) foi reservado exclusivamente para avaliação.

\subsection{Métrica de Execution Accuracy}
A métrica foi implementada como uma classe que herda de
\texttt{deepeval.metrics.BaseMetric}, com a lógica de avaliação no
método \texttt{measure}. O fluxo é: (i) extração robusta da consulta a
partir da saída bruta do modelo, priorizando blocos \texttt{```sql} e
recorrendo à primeira instrução \texttt{SELECT} no texto; (ii) conexão
somente-leitura ao banco SQLite correspondente ao \texttt{db\_id} do
caso de teste; (iii) execução da consulta gerada e da consulta gold sob
\texttt{try-except}, tratando erros de sintaxe como falha; (iv)
comparação dos conjuntos de resultados como multiconjuntos (insensível à
ordem das linhas), exceto quando a consulta gold contém
\texttt{ORDER BY}, caso em que a comparação preserva a ordem, seguindo a
semântica do avaliador oficial do Spider; (v) retorno de 1.0 para
resultados idênticos e 0.0 caso contrário. Esta métrica única foi
aplicada de forma idêntica ao baseline e aos modelos fine-tuned,
integrada a um teste \texttt{pytest} que consolida a acurácia agregada e
persiste o resultado por exemplo.

\subsection{Baseline}
O baseline foi obtido submetendo o modelo base ao dev split com um
prompt few-shot fixo contendo 3 exemplos (esquema + pergunta + SQL)
amostrados do training split com seed 42. Toda a decodificação de
avaliação foi determinística (greedy, \texttt{do\_sample=False}).

\subsection{Fine-tuning}
O fine-tuning empregou QLoRA: o modelo base quantizado em 4 bits (NF4,
double quantization, computação em bfloat16) com adaptadores LoRA. A
Tabela~\ref{tab:lora} documenta a configuração. Foram treinadas duas
configurações que diferem apenas na taxa de aprendizado, em uma ordem de
grandeza, para observar o impacto deste hiperparâmetro no trade-off.

\begin{table}[ht]
\centering
\caption{Configuração do LoRA e hiperparâmetros de treinamento.}
\label{tab:lora}
\begin{tabular}{lcc}
\toprule
Hiperparâmetro & Config.\ A & Config.\ B \\
\midrule
Rank ($r$) & 16 & 16 \\
Alpha & 32 & 32 \\
Dropout & 0.05 & 0.05 \\
Módulos alvo & \multicolumn{2}{c}{q,k,v,o,gate,up,down\_proj} \\
Taxa de aprendizado & $2\times10^{-4}$ & $2\times10^{-5}$ \\
Épocas & 2 & 2 \\
Batch efetivo & 16 & 16 \\
Comprimento máximo & 2048 & 2048 \\
Scheduler & cosine & cosine \\
Otimizador & \multicolumn{2}{c}{paged\_adamw\_8bit} \\
Seed & 42 & 42 \\
\bottomrule
\end{tabular}
\end{table}

\subsection{Avaliação de generalização (MMLU)}
A suíte de avaliação contém exatamente 150 questões: 50 de
\texttt{college\_computer\_science} (STEM), 50 de \texttt{philosophy}
(Humanidades) e 50 de \texttt{high\_school\_macroeconomics} (Ciências
Sociais), amostradas do test split com seed 42. A avaliação é 5-shot,
usando os 5 exemplos do dev split de cada disciplina como contexto
idêntico para o modelo base e os fine-tuned. A resposta é determinada
pela comparação dos logits dos tokens A--D na posição de resposta,
procedimento determinístico.

\section{Resultados}

\subsection{Tarefa-alvo (Spider)}

\begin{table}[ht]
\centering
\caption{Execution Accuracy no Spider dev split.}
\label{tab:spider}
\begin{tabular}{lc}
\toprule
Modelo & Execution Accuracy \\
\midrule
Baseline (few-shot) & [PREENCHER] \\
Fine-tuned A ($lr=2\times10^{-4}$) & [PREENCHER] \\
Fine-tuned B ($lr=2\times10^{-5}$) & [PREENCHER] \\
\bottomrule
\end{tabular}
\end{table}

[PREENCHER: 1--2 parágrafos descrevendo os números da
Tabela~\ref{tab:spider}: magnitude do ganho, diferença entre as duas
configurações.]

\subsection{Análise de erros}
[PREENCHER: examine 2--3 exemplos do arquivo
\texttt{results/eval\_predictions\_lora\_*.json} em que o modelo
fine-tuned falhou. Categorias típicas a procurar: erro de schema linking
(coluna/tabela inexistente), JOIN incorreto ou ausente, agregação errada
(COUNT vs.\ SUM), condição de filtro com valor literal errado. Mostre a
pergunta, a SQL gerada e a SQL gold de cada caso.]

\subsection{Regressão de capacidade (MMLU)}

\begin{table}[ht]
\centering
\caption{Acurácia 5-shot no MMLU (150 questões) e variação percentual em
relação ao baseline.}
\label{tab:mmlu}
\begin{tabular}{lccc}
\toprule
Categoria & Base & FT-A ($\Delta\%$) & FT-B ($\Delta\%$) \\
\midrule
STEM & [..] & [..] ([..]) & [..] ([..]) \\
Humanidades & [..] & [..] ([..]) & [..] ([..]) \\
C.\ Sociais & [..] & [..] ([..]) & [..] ([..]) \\
\midrule
Agregado & [..] & [..] ([..]) & [..] ([..]) \\
\bottomrule
\end{tabular}
\end{table}

[PREENCHER: parágrafo descrevendo a variação agregada e se a regressão
foi heterogênea entre categorias.]

\section{Discussão}

\subsection{O ganho justifica a perda?}
[PREENCHER conforme os números: compare a magnitude do ganho em
Execution Accuracy com a magnitude da perda no MMLU. Argumente em termos
do caso de uso: para um sistema dedicado a NL2SQL, perda moderada de
conhecimento geral tende a ser aceitável; para um assistente de
propósito geral, não.]

\subsection{Fatores que influenciam o trade-off}
[PREENCHER: relacione a diferença entre as configurações A e B. Hipótese
esperada: taxa de aprendizado maior tende a produzir maior ganho na
tarefa-alvo e maior regressão no MMLU; LoRA, por congelar os pesos base
e treinar apenas adaptadores de baixo rank, tende a mitigar o
esquecimento em comparação ao full fine-tuning. Discuta também o papel
do número de épocas e do rank.]

\subsection{Implicações práticas}
[PREENCHER: implicações para LLMs comerciais especializados ---
e.g., manutenção de suítes de regressão de capacidade como parte do
ciclo de release, uso de adaptadores intercambiáveis sobre um modelo
base congelado, mistura de dados gerais no conjunto de fine-tuning
(replay) para mitigar esquecimento.]

\subsection{Contaminação de dados}
Tanto o Spider quanto o MMLU são benchmarks públicos anteriores ao corte
de treinamento do modelo base, de modo que é plausível que parte de seus
exemplos tenha sido vista durante o pré-treino. Isso afeta a
interpretação em dois sentidos: o baseline no Spider pode estar inflado
(reduzindo o ganho aparente do fine-tuning) e a acurácia no MMLU pode
refletir memorização em vez de capacidade, o que tornaria a regressão
medida um limite inferior do esquecimento real. Como as comparações são
sempre relativas entre o mesmo modelo antes e depois do fine-tuning, sob
prompts e decodificação idênticos, as conclusões sobre a direção do
trade-off permanecem válidas, ainda que as magnitudes absolutas devam
ser lidas com cautela.

\section{Conclusão}
[PREENCHER: 1 parágrafo resumindo os achados quantitativos e a resposta
à pergunta central do trabalho.]

\begin{thebibliography}{9}
\bibitem{spider} T. Yu et al., ``Spider: A Large-Scale Human-Labeled
Dataset for Complex and Cross-Domain Semantic Parsing and Text-to-SQL
Task,'' EMNLP, 2018.
\bibitem{mmlu} D. Hendrycks et al., ``Measuring Massive Multitask
Language Understanding,'' ICLR, 2021.
\bibitem{lora} E. Hu et al., ``LoRA: Low-Rank Adaptation of Large
Language Models,'' ICLR, 2022.
\bibitem{qlora} T. Dettmers et al., ``QLoRA: Efficient Finetuning of
Quantized LLMs,'' NeurIPS, 2023.
\bibitem{qwen} Qwen Team, ``Qwen2.5 Technical Report,'' arXiv:2412.15115,
2024.
\end{thebibliography}

\end{document}
